% META: {"id": "TCOMPL-MF-CR-011", "chapitre": "TCOMPL-MODELES-FONCTION", "type_objet": "cours", "capacites_codes": ["C4", "C5"], "sources_inspiration": [], "mode_creation": "ex_nihilo", "status": "generated"}

\section{Convexité d'une fonction}

% ============================================================
% STRATE 1 — Reconnaitre la convexite graphiquement
% ============================================================

\definition[1]{
Une fonction $f$ dérivable sur un intervalle $I$ est \textbf{convexe} sur $I$ si sa courbe représentative est toujours située \textbf{au-dessus} de chacune de ses tangentes. Elle est \textbf{concave} si sa courbe est toujours \textbf{en dessous} de chacune de ses tangentes. Un \textbf{point d'inflexion} est un point où la courbe traverse sa tangente, c'est-à-dire où la convexité change.
}

\margeAppui{
Graphiquement : une courbe convexe \og se creuse vers le haut \fg{} (comme un bol), une courbe concave \og se creuse vers le bas \fg{} (comme un dôme).
}

% ============================================================
% STRATE 2 — Caracterisation par la derivee seconde
% ============================================================

\propriete[Caractérisation par la dérivée seconde]{
Soit $f$ deux fois dérivable sur un intervalle $I$. Alors :
\begin{itemize}
  \item $f$ est convexe sur $I$ si et seulement si $f''(x) \geq 0$ pour tout $x \in I$ ;
  \item $f$ est concave sur $I$ si et seulement si $f''(x) \leq 0$ pour tout $x \in I$ ;
  \item un point d'inflexion correspond à un point où $f''$ s'annule en \textbf{changeant de signe}.
\end{itemize}
}

\exemple{
Reprenons $f(x) = x + \dfrac{9}{x}$ sur $]0,+\infty[$. On calcule $f''(x) = \dfrac{18}{x^3}$. Pour tout $x>0$, $f''(x)>0$ : $f$ est \textbf{convexe} sur $]0,+\infty[$, sans point d'inflexion.
}

% BEGIN-VERIFY
% from sympy import symbols, diff
% x = symbols('x', positive=True)
% f = x + 9/x
% fpp = diff(f, x, 2)
% assert fpp == 18/x**3
% END-VERIFY

\exemple{
Soit $g(x) = x^3 - 3x$ sur $\mathbb{R}$. On calcule $g''(x) = 6x$, qui s'annule en $x=0$ en changeant de signe (négatif pour $x<0$, positif pour $x>0$) : $g$ admet un \textbf{point d'inflexion} en $x=0$, elle est concave sur $]-\infty,0]$ et convexe sur $[0,+\infty[$.
}

% BEGIN-VERIFY
% from sympy import symbols, diff, solve, Eq
% x = symbols('x')
% g = x**3 - 3*x
% gpp = diff(g, x, 2)
% assert gpp == 6*x
% assert solve(Eq(gpp, 0), x) == [0]
% assert gpp.subs(x, -1) < 0
% assert gpp.subs(x, 1) > 0
% END-VERIFY

\erreurFrequente{
\textbf{Confondre $f'' \geq 0$ et $f' \geq 0$.} Le signe de $f'$ donne le sens de \textbf{variation} de $f$ (croissante/décroissante) ; le signe de $f''$ donne la \textbf{convexité} de $f$ (courbe au-dessus/en dessous des tangentes). Une fonction peut être croissante et concave, croissante et convexe, décroissante et concave, ou décroissante et convexe : les deux notions sont indépendantes.
}
